\section{Three years worked end to end}
\label{sec:years}

The rules of Sections~\ref{sec:ot} and~\ref{sec:cycles} are worth nothing unless they produce the right answer.
Three consecutive liturgical years are therefore carried through in full below, each computed from
the norms alone and then checked, day by day at the points where the computation could fail, against
the \emph{Liturgical Calendar for the Dioceses of the United States of America} published by the
Secretariat of Divine Worship for that year. Two of the three are thirty-three-week years with an
omitted week; the third is a thirty-four-week year, included precisely so that the rule can be seen
not to fire.

All three are worked for the United States, where the Epiphany is assigned to the Sunday from 2 to 8
January and the Body and Blood of Christ to the Sunday after Trinity. Where the Ascension differs by
province, both branches are given.

\subsection{Liturgical year 2026: thirty-three weeks, week VII omitted}

\begin{ruletable}
\textbf{Boundaries} & Begins at Evening Prayer I of the First Sunday of Advent, 30 November 2025;
ends Saturday 28 November 2026. $Y$ = 2026.\\
\textbf{Cycles} & 2026 leaves remainder 1 on division by 3, so Sunday Year~A; 2026 is even, so weekday Cycle~II. The
Secretariat states Year~A from 30 November 2025 to 22 November 2026, and Cycle~2 from 12 January to
17 February and from 25 May to 28 November 2026.\\
\textbf{Christmas Time} & Nativity Thursday 25 December 2025; Holy Family Sunday 28 December; Mary,
Mother of God, Thursday 1 January 2026; Epiphany Sunday 4 January; Baptism of the Lord Sunday 11
January. There is no Second Sunday after the Nativity: the only Sunday available to it, 4 January,
is the Epiphany.\\
\textbf{Ordinary Time begins} & Monday 12 January, the Monday after the Sunday following 6 January
(\loc{nualc}~44). Week~I runs 12--17 January and has no numbered Sunday.\\
\textbf{Pre-Lenten Sundays} & Second Sunday 18 January through Sixth Sunday 15 February. Ash
Wednesday is 18 February, so $L$ = 6.\\
\textbf{The truncated week} & The readings of week VI are used on Monday 16 and Tuesday 17 February
and then suspended (\loc{olm}~104.2). The Wednesday-to-Saturday sets of week VI --- Lectionary
337--340 --- are not read.\\
\textbf{Easter cycle} & Easter 5 April; Ascension Thursday 14 May in the provinces of Boston,
Hartford, New York, Omaha and Philadelphia, and Sunday 17 May elsewhere; Pentecost Sunday 24 May.\\
\textbf{Resumption} & The Sunday before Advent is 22 November; 24 May to 22 November is
26 whole weeks, so $R$ = 34 $-$ 26 = 8. Ordinary Time resumes on Monday 25 May in week VIII. The Secretariat
labels 25 May ``Eighth Week in Ordinary Time'' and assigns the obligatory memorial of the Blessed
Virgin Mary, Mother of the Church, to that Monday.\\
\textbf{The omission} & $R$ = $L$ + 2, so the year has thirty-three weeks and week VII never occurs.
The Seventh Sunday in Ordinary Time does not appear in the calendar; its Missal formulary is unused;
the weekday sets 341--346 are not read. With the four sets lost from week VI, ten weekday sets go
unread in 2026.\\
\textbf{The Sundays after Pentecost} & Pentecost is the Sunday of week VIII, the Most Holy Trinity
(31 May) of week IX, the Body and Blood of Christ (7 June) of week X. The first numbered Sunday to
reappear is therefore the Eleventh, on 14 June --- which is exactly what the Secretariat prints.\\
\textbf{Close of the year} & Christ the King Sunday 22 November, the Sunday of week XXXIV; the week's
weekdays run to Saturday 28 November; the First Sunday of Advent 2027 is 29 November 2026.\\
\textbf{Independent corroboration} & The Secretariat's Liturgy of the Hours table for 2026 assigns
volume III to ``Weeks 1 to 6, Ordinary Time'' from 12 January to 17 February and again to ``Weeks 8
to 17, Ordinary Time'' from 25 May to 1 August, with volume IV taking ``Weeks 18 to 34'' from 2
August to 28 November. The gap between 6 and 8 is the omitted week, recorded in an entirely
different apparatus of the same document.\\
\end{ruletable}

\subsection{Liturgical year 2027: thirty-three weeks, week VI omitted}

\begin{ruletable}
\textbf{Boundaries} & 29 November 2026 to Saturday 27 November 2027. $Y$ = 2027.\\
\textbf{Cycles} & 2027 leaves remainder 2 on division by 3, so Sunday Year~B; 2027 odd, weekday Cycle~I. The Secretariat gives
Year~B from 29 November 2026 to 21 November 2027 and Cycle~1 from 11 January to 9 February and from
17 May to 27 November 2027.\\
\textbf{Christmas Time} & Epiphany Sunday 3 January 2027; Baptism of the Lord Sunday 10 January.\\
\textbf{Ordinary Time begins} & Monday 11 January; week~I runs 11--16 January.\\
\textbf{Pre-Lenten Sundays} & Second Sunday 17 January, Third 24 January, Fourth 31 January, Fifth 7
February. Ash Wednesday 10 February, so $L$ = 5 --- the exact case worked in \loc{olm}~104 footnote
116.\\
\textbf{Easter cycle} & Easter 28 March; Ascension Thursday 6 May or Sunday 9 May by province;
Pentecost Sunday 16 May.\\
\textbf{Resumption} & Sunday before Advent 21 November; 16 May to 21 November is 27 whole weeks, so
$R$ = 7. Ordinary Time resumes Monday 17 May in week VII, and the Secretariat's entry for that Monday
reads ``Seventh Week in Ordinary Time.''\\
\textbf{The omission} & $R$ = $L$ + 2: thirty-three weeks, week VI omitted. Footnote 116 predicts this
result in the same words --- five weeks before Lent, the sixth omitted, begin from the seventh.\\
\textbf{The Sundays after Pentecost} & Pentecost is the Sunday of week VII, Trinity (23 May) of week
VIII, the Body and Blood of Christ (30 May) of week IX. The first numbered Sunday is the Tenth, 6
June, as printed.\\
\textbf{Close of the year} & Christ the King 21 November; last day Saturday 27 November; Advent
begins 28 November 2027.\\
\end{ruletable}

\subsection{Liturgical year 2028: thirty-four weeks, nothing omitted}

\begin{ruletable}
\textbf{Boundaries} & 28 November 2027 to Saturday 2 December 2028. $Y$ = 2028.\\
\textbf{Cycles} & 2028 divides exactly by 3, so Sunday Year~C; 2028 even, weekday Cycle~II. The Secretariat
gives Year~C from 28 November 2027 to 26 November 2028 and Cycle~II from 10 January to 29 February
and from 5 June to 2 December 2028.\\
\textbf{Christmas Time} & Epiphany Sunday 2 January 2028; Baptism of the Lord Sunday 9 January.\\
\textbf{Ordinary Time begins} & Monday 10 January; week~I runs 10--15 January.\\
\textbf{Pre-Lenten Sundays} & Second Sunday 16 January through Eighth Sunday 27 February. Ash
Wednesday 1 March, so $L$ = 8. Because Ash Wednesday is a Wednesday of week VIII, the readings of
that week are still suspended after Tuesday 29 February.\\
\textbf{Easter cycle} & Easter 16 April; Ascension Thursday 25 May or Sunday 28 May by province;
Pentecost Sunday 4 June.\\
\textbf{Resumption} & Sunday before Advent 26 November; 4 June to 26 November is 25 whole weeks, so
$R$ = 9. Ordinary Time resumes Monday 5 June in week IX, and the Secretariat prints ``Ninth Week in
Ordinary Time'' for the Monday after Pentecost.\\
\textbf{No omission} & $R$ = $L$ + 1: thirty-four weeks. Every Missal formulary and every weekday
Lectionary set of Ordinary Time is available in 2028, except the four sets of week VIII lost to Ash
Wednesday. This is the case of \loc{olm}~104 footnote 115.\\
\textbf{The Sundays after Pentecost} & Pentecost is the Sunday of week IX, Trinity (11 June) of week
X, the Body and Blood of Christ (18 June) of week XI. The first numbered Sunday is the Twelfth, 25
June, as printed.\\
\textbf{Close of the year} & Christ the King 26 November; last day Saturday 2 December; Advent begins
3 December 2028.\\
\end{ruletable}

\subsection{The three years side by side}

\begin{comparetable}
Last pre-Lent week $L$ & 6 / 5 / 8\\
Resumed week $R$ & 8 / 7 / 9\\
$R - L$ & 2 / 2 / 1\\
Weeks in Ordinary Time & 33 / 33 / 34\\
Week omitted & VII / VI / none\\
First numbered Sunday after Pentecost & XI, 14 June / X, 6 June / XII, 25 June\\
Sunday cycle & A / B / C\\
Weekday cycle & II / I / II\\
\end{comparetable}

The three Sunday letters and the two weekday numerals in this window are themselves a demonstration
of \loc{olm}~65: over three consecutive years the Sunday cycle completes exactly one revolution while
the weekday cycle completes one and a half, and the pairing A--II, B--I, C--II does not repeat until
2032, 2033 and 2034.

\subsection{A longer computed table}

The table below is computed from the rules of this reference for twelve liturgical years, on the
United States assignments. The rows for 2026, 2027 and 2028 are the ones checked above against the
competent territorial calendar; the remaining rows are computation only and carry no official
witness. Dates are of the civil year named in the first column, except the First Sunday of Advent,
which opens the \emph{following} liturgical year.

\begin{yeartable}
2024 & 8 Jan (Mon) & 14 Feb & 31 Mar & 19 May & 1 Dec & 6 & 7 & 34 & --- & B & II\\
2025 & 12 Jan & 5 Mar & 20 Apr & 8 Jun & 30 Nov & 8 & 10 & 33 & IX & C & I\\
2026 & 11 Jan & 18 Feb & 5 Apr & 24 May & 29 Nov & 6 & 8 & 33 & VII & A & II\\
2027 & 10 Jan & 10 Feb & 28 Mar & 16 May & 28 Nov & 5 & 7 & 33 & VI & B & I\\
2028 & 9 Jan & 1 Mar & 16 Apr & 4 Jun & 3 Dec & 8 & 9 & 34 & --- & C & II\\
2029 & 8 Jan (Mon) & 14 Feb & 1 Apr & 20 May & 2 Dec & 6 & 7 & 34 & --- & A & I\\
2030 & 13 Jan & 6 Mar & 21 Apr & 9 Jun & 1 Dec & 8 & 10 & 33 & IX & B & II\\
2031 & 12 Jan & 26 Feb & 13 Apr & 1 Jun & 30 Nov & 7 & 9 & 33 & VIII & C & I\\
2032 & 11 Jan & 11 Feb & 28 Mar & 16 May & 28 Nov & 5 & 7 & 33 & VI & A & II\\
2033 & 9 Jan & 2 Mar & 17 Apr & 5 Jun & 27 Nov & 8 & 10 & 33 & IX & B & I\\
2034 & 9 Jan (Mon) & 22 Feb & 9 Apr & 28 May & 3 Dec & 7 & 8 & 34 & --- & C & II\\
2035 & 8 Jan (Mon) & 7 Feb & 25 Mar & 13 May & 2 Dec & 5 & 6 & 34 & --- & A & I\\
\end{yeartable}

The Baptism of the Lord is marked \emph{(Mon)} in the years in which the Epiphany, assigned to
Sunday, falls on 7 or 8 January and displaces it under the Missal's rubric. In those years the
readings of week~I begin on the Tuesday, not the Monday, under \loc{olm}~104.1, and week~I has five
weekdays instead of six.

\subsection{The Missal's own table as a control}

Section~\ref{sec:tabella} reported the recomputation of the \latin{Tabella temporaria} of the 2002
\latin{editio typica tertia}. It bears restating here as verification rather than as bibliography:
for the twenty-four civil years 2000--2023 the promulgated table records, in its own columns, the
last pre-Lent week number and the first resumed week number for each year, and in all twenty-four
those numbers agree with the rule published in Section~\ref{sec:ot}. Together with the three years
checked against the United States calendar, the arithmetic of this reference has been tested against
an official witness in twenty-seven distinct years, and the two witnesses are independent of each
other.
